% =================================================================
%  A Post-Quantum Crypto Agility Framework: Runtime Algorithm
%  Negotiation, Quantum Threat Scoring, and Hybrid Key Exchange
%
%  Target: IACR Communications in Cryptology (CiC)
% =================================================================
\documentclass[conference,10pt]{IEEEtran}

% ── packages ────────────────────────────────────────────────────
\usepackage[utf8]{inputenc}
\usepackage[T1]{fontenc}
\usepackage{amsmath,amssymb,amsfonts}
\usepackage{algorithmic}
\usepackage{algorithm}
\usepackage{graphicx}
\usepackage{textcomp}
\usepackage{booktabs}
\usepackage{multirow}
\usepackage{xcolor}
\usepackage{url}
\usepackage{hyperref}
\usepackage{cite}
\usepackage{balance}
\usepackage{tabularx}
\usepackage{enumitem}

\hypersetup{
  colorlinks=true,
  linkcolor=blue!70!black,
  citecolor=blue!70!black,
  urlcolor=blue!70!black
}

% ── macros ──────────────────────────────────────────────────────
\newcommand{\qrs}{\textrm{QRS}}
\newcommand{\avs}{\textrm{AVS}}
\newcommand{\thr}{\textrm{THR}}
\newcommand{\dss}{\textrm{DSS}}
\newcommand{\mgs}{\textrm{MGS}}
\newcommand{\es}{\textrm{ES}}
\newcommand{\mlkem}{\textrm{ML-KEM}}
\newcommand{\mldsa}{\textrm{ML-DSA}}
\newcommand{\slhdsa}{\textrm{SLH-DSA}}
\newcommand{\system}{\textsc{QBridge}}

% =================================================================
\begin{document}

\title{A Post-Quantum Crypto Agility Framework:\\Runtime Algorithm Negotiation, Quantum Threat Scoring, and Hybrid Key Exchange}

\author{
  \IEEEauthorblockN{Prabakaran Kannan}
  \IEEEauthorblockA{Research Scholar, Department of Computer Science and Engineering\\
  National Institute of Technology Puducherry\\
Email: cs23d2005@nitpy.ac.in}
}

\maketitle

% =================================================================
%  ABSTRACT
% =================================================================
\begin{abstract}
The imminent arrival of cryptographically relevant quantum computers (CRQCs)
demands that deployed systems migrate away from RSA and elliptic-curve
cryptography toward post-quantum alternatives standardized by NIST.
However, enterprise environments encompass thousands of heterogeneous
cryptographic assets whose simultaneous migration is infeasible, creating
a protracted transition period during which classical, hybrid, and
post-quantum algorithms must coexist. We present \system{}, an integrated
framework addressing three critical facets of this transition.
First, we introduce a \emph{Crypto Agility Negotiation Protocol} that
enables runtime algorithm selection between peers through capability
advertisement, policy-constrained intersection, and zero-downtime
rotation when algorithms are deprecated or compromised. The protocol
supports 28~algorithms across six families governed by four organizational
policies aligned with NIST~SP~800-227, CNSA~2.0, and ETSI~TS~103\,744.
Second, we describe a \emph{Quantum Threat Scoring Engine} that computes
a composite Quantum Risk Score (QRS) in $[0,100]$ by weighting five
factors---algorithm vulnerability, time-horizon risk via Mosca's
inequality, data sensitivity, migration gap, and exposure surface---to
prioritize migration across an entire cryptographic asset portfolio.
Third, we specify three \emph{Hybrid Key Exchange} constructions
(X25519\nobreakdash-ML\nobreakdash-KEM\nobreakdash-768,
P384\nobreakdash-ML\nobreakdash-KEM\nobreakdash-1024,
X25519\nobreakdash-ML\nobreakdash-KEM\nobreakdash-512) that combine
classical ECDH with Module-LWE-based key encapsulation, deriving the
shared secret through HKDF-SHA256 with domain separation, suitable for
TLS~1.3 integration. We evaluate \system{} on realistic workloads,
demonstrating negotiation latency under 2\,ms, scoring throughput
exceeding 50\,000 assets per second, and hybrid KEM overhead within
15\% of standalone ML-KEM, making the framework practical for
large-scale deployment.
\end{abstract}

\begin{IEEEkeywords}
post-quantum cryptography, crypto agility, hybrid key exchange,
ML-KEM, quantum risk assessment, TLS 1.3, NIST PQC standards
\end{IEEEkeywords}

% =================================================================
%  I. INTRODUCTION
% =================================================================
\section{Introduction}\label{sec:intro}

Shor's algorithm~\cite{shor1997} renders RSA, finite-field
Diffie--Hellman, and elliptic-curve cryptography insecure against a
sufficiently large quantum computer. Although the exact timeline for a
cryptographically relevant quantum computer (CRQC) remains uncertain,
conservative estimates place this event within the 2030--2040
window~\cite{mosca2018}. The \emph{harvest-now, decrypt-later} (HNDL)
threat model makes the situation more acute: adversaries can record
ciphertexts today and decrypt them once a CRQC becomes
available~\cite{mavroeidis2018}. For data with long confidentiality
requirements---financial records, medical data, classified
communications---migration to quantum-resistant algorithms must begin
\emph{now}.

NIST finalized its first post-quantum standards in 2024: FIPS~203
(\mlkem{})~\cite{nist-fips203} for key encapsulation, FIPS~204
(\mldsa{})~\cite{nist-fips204} for digital signatures, and FIPS~205
(\slhdsa{})~\cite{nist-fips205} for stateless hash-based signatures.
Supplementary guidance in SP~800-227~\cite{nist-sp800-227} provides a
migration roadmap, while the NSA's CNSA~2.0 suite~\cite{cnsa2} imposes
stricter timelines for national security systems. Despite these
standards, the practical challenge of migrating a large enterprise from
legacy cryptography to PQC remains daunting.

Three obstacles stand out.
\emph{(i)~Heterogeneity.}  A typical enterprise deploys dozens of
distinct cryptographic algorithms across TLS termination, code signing,
database encryption, VPN tunnels, and API authentication. No single
algorithm or policy applies uniformly.
\emph{(ii)~Risk prioritization.}  With limited migration resources,
organizations must decide \emph{which} assets to migrate first.
Intuitive approaches (``migrate everything using RSA'') ignore data
sensitivity, exposure surface, and the HNDL timeline.
\emph{(iii)~Transition coexistence.}  During the multi-year migration
window, hybrid constructions that combine classical and PQC algorithms
are necessary both for defense-in-depth and for interoperability with
peers that have not yet migrated~\cite{bindel2019}.

\subsection{Contributions}

We present \system{}, an integrated crypto agility framework that
addresses all three obstacles through the following contributions:

\begin{enumerate}[leftmargin=*]
\item \textbf{Crypto Agility Negotiation Protocol.}  A runtime
  algorithm negotiation mechanism that advertises cryptographic
  capabilities, computes policy-constrained intersections, and supports
  zero-downtime algorithm rotation (Section~\ref{sec:agility}).
  The protocol manages 28 algorithms across six families under four
  configurable policies.

\item \textbf{Quantum Threat Scoring Engine.}  A composite risk model
  that assigns a Quantum Risk Score (QRS) to every cryptographic
  asset in an organization's portfolio, enabling quantitative migration
  prioritization based on Mosca's inequality and five weighted factors
  (Section~\ref{sec:scoring}).

\item \textbf{Hybrid Key Exchange Constructions.}  Three hybrid KEM
  variants for TLS~1.3---X25519-\mlkem{}-768, P384-\mlkem{}-1024, and
  X25519-\mlkem{}-512---with HKDF-based shared-secret combination
  and domain separation (Section~\ref{sec:hybrid}).

\item \textbf{Implementation and Evaluation.}  A production-quality
  Python implementation with performance benchmarks demonstrating
  practical deployment characteristics (Section~\ref{sec:eval}).
\end{enumerate}

The remainder of the paper is organized as follows.
Section~\ref{sec:background} covers background and preliminaries.
Sections~\ref{sec:agility}--\ref{sec:hybrid} detail our three
constructions. Section~\ref{sec:eval} presents implementation results.
Section~\ref{sec:security} provides a security analysis.
Section~\ref{sec:related} surveys related work, and
Section~\ref{sec:conclusion} concludes.

% =================================================================
%  II. BACKGROUND AND PRELIMINARIES
% =================================================================
\section{Background and Preliminaries}\label{sec:background}

\subsection{Lattice-Based Cryptography}

The NIST PQC standards are built primarily on the hardness of the
\emph{Module Learning with Errors} (Module-LWE) and \emph{Module Short
Integer Solution} (Module-SIS) problems~\cite{langlois2015}.

\begin{definition}[Module-LWE]
Given a rank-$k$ module over the polynomial ring
$R_q = \mathbb{Z}_q[X]/(X^n+1)$, the Module-LWE problem asks to
distinguish $(A, A \cdot s + e)$ from $(A, u)$ where
$A \xleftarrow{\$} R_q^{k \times k}$,
$s \xleftarrow{} \chi_s^k$,
$e \xleftarrow{} \chi_e^k$, and
$u \xleftarrow{\$} R_q^k$, with $\chi_s, \chi_e$ being centered
binomial distributions with small parameter $\eta$.
\end{definition}

ML-KEM (FIPS~203) instantiates an IND-CCA2 key encapsulation mechanism
from Module-LWE via the Fujisaki--Okamoto transform.  Three parameter
sets---ML-KEM-512 (NIST Level~1), ML-KEM-768 (Level~3), and ML-KEM-1024
(Level~5)---provide increasing security margins at the cost of larger
keys and ciphertexts~\cite{nist-fips203}.

ML-DSA (FIPS~204) is a digital signature scheme based on
Module-LWE/Module-SIS with parameter sets ML-DSA-44, ML-DSA-65, and
ML-DSA-87~\cite{nist-fips204}. SLH-DSA (FIPS~205) provides a
conservative alternative based solely on hash-function
security~\cite{nist-fips205}.

\subsection{Quantum Threat Model}

We adopt the standard \emph{harvest-now, decrypt-later} (HNDL) threat
model~\cite{mavroeidis2018}. An adversary with classical computational
power records ciphertexts and authenticated protocol transcripts today.
At a future time $t_{\text{CRQC}}$, the adversary gains access to a CRQC
capable of running Shor's algorithm~\cite{shor1997} and recovers
plaintext from all recorded sessions using vulnerable key exchange
algorithms.

Symmetric algorithms are affected by Grover's algorithm~\cite{grover1996},
which provides a quadratic speedup for brute-force search, effectively
halving the security level (e.g., AES-128 provides only 64-bit security
against a quantum adversary, while AES-256 retains 128-bit security).

\subsection{Mosca's Inequality}

Mosca's theorem~\cite{mosca2018} provides a decision criterion for
migration urgency:

\begin{equation}\label{eq:mosca}
  x + y > z \implies \text{migration is overdue}
\end{equation}

\noindent where $x$ is the number of years the data must remain
confidential, $y$ is the time (in years) required to complete migration,
and $z$ is the estimated time until a CRQC can break the algorithm. When
$x + y > z$, data encrypted today will still require protection after the
quantum threat materializes, and the migration cannot complete in time.

\subsection{Standards Landscape}

Our framework is aligned with four major standards:

\begin{itemize}[leftmargin=*]
\item \textbf{NIST SP 800-227}~\cite{nist-sp800-227}: Provides a
  phased transition roadmap for federal agencies to adopt PQC.
\item \textbf{CNSA 2.0}~\cite{cnsa2}: The NSA's algorithm suite
  mandating ML-KEM-1024, ML-DSA-87, and AES-256 for national security
  systems by 2035.
\item \textbf{ETSI TS 103 744}~\cite{etsi-ts103744}: European
  migration strategies including hybrid approaches.
\item \textbf{RFC 8446}~\cite{rfc8446}: TLS 1.3 cipher suite
  negotiation, which our agility protocol extends for PQC.
\end{itemize}

% =================================================================
%  III. CRYPTO AGILITY NEGOTIATION PROTOCOL
% =================================================================
\section{Crypto Agility Negotiation Protocol}\label{sec:agility}

Crypto agility---the ability to switch algorithms without redesigning
protocols---has been identified as a critical requirement for the PQC
transition~\cite{ott2019, housley2024, dhs-pqc-roadmap}. We formalize a
negotiation protocol that supports runtime algorithm selection, policy
enforcement, and zero-downtime rotation.

\subsection{Design Goals}

Our protocol satisfies five design goals:

\begin{enumerate}[leftmargin=*, label=\textbf{G\arabic*.}]
\item \textbf{Forward Compatibility:} Unknown future algorithms can be
  negotiated via OID or canonical name without protocol changes.
\item \textbf{Hybrid Negotiation:} Classical and PQC algorithms
  operate simultaneously during the transition period.
\item \textbf{Policy Enforcement:} Organizational policies constrain
  negotiation to approved algorithms with minimum security levels.
\item \textbf{Zero-Downtime Rotation:} Algorithms can be banned or
  upgraded without interrupting active sessions.
\item \textbf{Audit Trail:} Complete negotiation history is maintained
  for compliance reporting.
\end{enumerate}

\subsection{Algorithm Registry}

The protocol maintains a registry of 28 algorithms organized into six
families: KEM, signature, symmetric, hash, MAC, and KDF.
Table~\ref{tab:algorithm-registry} summarizes the key entries.

\begin{table}[t]
\centering
\caption{Selected entries from the algorithm registry.}
\label{tab:algorithm-registry}
\footnotesize
\begin{tabular}{@{}llrrrr@{}}
\toprule
\textbf{Algorithm} & \textbf{Family} & \textbf{Bits} & \textbf{Level} & \textbf{PK (B)} & \textbf{CT/Sig (B)} \\
\midrule
ML-KEM-512         & KEM  & 128 & 1 & 800   & 768   \\
ML-KEM-768         & KEM  & 192 & 3 & 1\,184 & 1\,088 \\
ML-KEM-1024        & KEM  & 256 & 5 & 1\,568 & 1\,568 \\
X25519-ML-KEM-768  & KEM  & 192 & 3 & 1\,216 & 1\,120 \\
P384-ML-KEM-1024   & KEM  & 256 & 5 & 1\,665 & 1\,665 \\
\midrule
ML-DSA-44          & Sig  & 128 & 2 & 1\,312 & 2\,420 \\
ML-DSA-65          & Sig  & 192 & 3 & 1\,952 & 3\,293 \\
ML-DSA-87          & Sig  & 256 & 5 & 2\,592 & 4\,595 \\
Falcon-512         & Sig  & 128 & 1 & 897   & 666   \\
SLH-DSA-128f       & Sig  & 128 & 1 & 32    & 17\,088 \\
\midrule
AES-256-GCM        & Sym  & 256 & -- & 32    & --    \\
SHA3-256            & Hash & 256 & -- & --    & 32    \\
HKDF-SHA3-256      & KDF  & 256 & -- & --    & --    \\
\bottomrule
\end{tabular}
\end{table}

Each algorithm is assigned a lifecycle status from the ordered set
$\{\textsc{Recommended}, \textsc{Acceptable}, \textsc{Legacy},
\textsc{Deprecated}, \textsc{Banned}\}$, and a security era from
$\{\textsc{Classical}, \textsc{Hybrid}, \textsc{Post-Quantum},
\textsc{CNSA-2.0}\}$.

\subsection{Organizational Policies}

Four pre-defined policies govern negotiation behavior:

\begin{description}[leftmargin=*]
\item[Transitional:] $\geq 128$-bit security, hybrid required,
  classical-only disallowed. Suitable for organizations beginning PQC
  migration.

\item[Post-Quantum:] $\geq 192$-bit security, NIST Level~$\geq 3$,
  classical algorithms banned. For environments ready for full PQC
  deployment.

\item[CNSA 2.0:] $\geq 256$-bit security, NIST Level~5, prefers
  ML-KEM-1024 and ML-DSA-87. Targets national security compliance.

\item[Legacy Compatible:] $\geq 128$-bit security, classical-only
  permitted. For backward compatibility with legacy peers.
\end{description}

\subsection{Negotiation Protocol}

The negotiation proceeds in three phases, formalized in
Algorithm~\ref{alg:negotiate}.

\begin{algorithm}[t]
\caption{Crypto Agility Negotiation}
\label{alg:negotiate}
\begin{algorithmic}[1]
\renewcommand{\algorithmicrequire}{\textbf{Input:}}
\renewcommand{\algorithmicensure}{\textbf{Output:}}
\REQUIRE Local policy $\pi$, local capability $C_L$, remote capability $C_R$
\ENSURE Negotiated suite $\mathcal{S}$ or $\perp$
\STATE \textbf{Phase 1: Capability Advertisement}
\STATE $C_L \leftarrow \textsc{FilterByPolicy}(\pi, \text{Registry})$
\STATE Send $C_L$ to remote peer
\STATE Receive $C_R$ from remote peer
\STATE
\STATE \textbf{Phase 2: Policy-Constrained Selection}
\FOR{each family $f \in \{\text{KEM}, \text{Sig}, \text{Sym}, \text{Hash}, \text{KDF}\}$}
  \STATE $I_f \leftarrow C_L[f] \cap C_R[f]$
  \IF{$I_f = \emptyset$}
    \RETURN $\perp$ \COMMENT{Negotiation failure}
  \ENDIF
  \STATE $a_f \leftarrow$ first element of $I_f$ by local preference
  \IF{$a_f.\text{status} = \textsc{Banned}$ \OR $a_f.\text{bits} < \pi.\text{min\_bits}$}
    \STATE Remove $a_f$ from $I_f$; retry selection
  \ENDIF
\ENDFOR
\STATE
\STATE \textbf{Phase 3: Suite Construction}
\STATE $\text{era} \leftarrow \textsc{DetermineEra}(a_\text{KEM}, a_\text{Sig}, a_\text{Sym}, a_\text{Hash})$
\STATE $\text{id} \leftarrow \text{SHA3-256}(a_\text{KEM} \| a_\text{Sig} \| a_\text{Sym} \| a_\text{Hash} \| r)$, $r \xleftarrow{\$} \{0,1\}^{64}$
\STATE $\mathcal{S} \leftarrow (a_\text{KEM}, a_\text{Sig}, a_\text{Sym}, a_\text{Hash}, a_\text{KDF}, \text{era}, \text{id})$
\RETURN $\mathcal{S}$
\end{algorithmic}
\end{algorithm}

The \textsc{DetermineEra} function classifies the negotiated suite:
if all selected algorithms belong to the post-quantum or CNSA-2.0 era,
the suite era is set to \textsc{Post-Quantum}; if at least one algorithm
is post-quantum while others are classical, the era is \textsc{Hybrid};
otherwise the era is \textsc{Classical}.

\subsection{Zero-Downtime Rotation}

When a vulnerability is discovered in an active algorithm, the framework
supports in-place renegotiation without session teardown:

\begin{algorithm}[t]
\caption{Zero-Downtime Algorithm Rotation}
\label{alg:rotate}
\begin{algorithmic}[1]
\REQUIRE Current suite $\mathcal{S}$, algorithm to ban $a_{\text{ban}}$, optional target era $e_{\text{new}}$
\ENSURE Updated policy $\pi'$
\STATE $B \leftarrow \pi.\text{banned} \cup \{a_{\text{ban}}\}$
\STATE $\pi' \leftarrow \pi$ with $\text{banned} \leftarrow B$
\IF{$e_{\text{new}} \neq \textsc{null}$}
  \STATE $\pi'.\text{required\_era} \leftarrow e_{\text{new}}$
\ENDIF
\STATE Trigger renegotiation with $\pi'$
\RETURN $\pi'$
\end{algorithmic}
\end{algorithm}

The rotation is atomic with respect to the session: the current suite
remains active until a new suite is successfully negotiated, ensuring no
interruption in cryptographic protection.

\subsection{Compliance Auditing}

Every negotiation event is appended to an immutable history with the
following fields: suite identifier, selected algorithms, security era,
initiator and responder identities, timestamp, and validity period. This
history supports compliance reporting against NIST~SP~800-227 and
CNSA~2.0 timelines.

% =================================================================
%  IV. QUANTUM THREAT SCORING ENGINE
% =================================================================
\section{Quantum Threat Scoring Engine}\label{sec:scoring}

Risk-based prioritization is essential when migration resources are
limited~\cite{joseph2022}. We formalize a scoring model that assigns
each cryptographic asset a Quantum Risk Score (QRS), enabling
quantitative comparison across an organization's portfolio.

\subsection{Composite Risk Model}

The QRS for a cryptographic asset $a$ is defined as:

\begin{equation}\label{eq:qrs}
  \qrs(a) = \sum_{i=1}^{5} w_i \cdot f_i(a)
\end{equation}

\noindent where each factor $f_i \in [0, 100]$ and the weights $w_i$
satisfy $\sum w_i = 1$. Table~\ref{tab:weights} lists the default
weighting.

\begin{table}[t]
\centering
\caption{QRS factor weights.}
\label{tab:weights}
\footnotesize
\begin{tabular}{@{}llc@{}}
\toprule
\textbf{Factor} & \textbf{Symbol} & \textbf{Weight $w_i$} \\
\midrule
Algorithm Vulnerability Score & $\avs$ & 0.30 \\
Time Horizon Risk             & $\thr$ & 0.25 \\
Data Sensitivity Score        & $\dss$ & 0.20 \\
Migration Gap Score           & $\mgs$ & 0.15 \\
Exposure Surface              & $\es$  & 0.10 \\
\bottomrule
\end{tabular}
\end{table}

\subsection{Factor Definitions}

\paragraph{Algorithm Vulnerability Score ($\avs$)}
Each algorithm is assigned a base vulnerability score reflecting its
susceptibility to quantum attack. Algorithms vulnerable to Shor's
algorithm receive scores of $75$--$90$ (e.g., ECDSA-P256: $90$,
RSA-2048: $85$). Algorithms affected only by Grover's attack score
$5$--$30$ (e.g., AES-128: $30$, AES-256: $5$). NIST PQC algorithms
score $1$--$5$ (e.g., ML-KEM-768: $2$, ML-KEM-1024: $1$).
Table~\ref{tab:avs} presents the complete scoring matrix.

\begin{table}[t]
\centering
\caption{Algorithm Vulnerability Scores.}
\label{tab:avs}
\footnotesize
\begin{tabular}{@{}lccl@{}}
\toprule
\textbf{Algorithm} & \textbf{AVS} & \textbf{Shor} & \textbf{Quantum Impact} \\
\midrule
ECDSA-P256       & 90  & Yes & Fully broken \\
ECDH-P256        & 90  & Yes & Fully broken \\
RSA-2048         & 85  & Yes & Fully broken \\
RSA-3072         & 80  & Yes & Fully broken \\
ECDSA-P384       & 85  & Yes & Fully broken \\
DH-2048          & 85  & Yes & Fully broken \\
\midrule
AES-128          & 30  & No  & Security halved (64-bit) \\
SHA-256          & 15  & No  & Collision reduced \\
AES-256          & 5   & No  & Security halved (128-bit) \\
SHA3-256         & 5   & No  & Minimal impact \\
\midrule
ML-KEM-768       & 2   & No  & Quantum-resistant \\
ML-KEM-1024      & 1   & No  & Quantum-resistant \\
ML-DSA-65        & 2   & No  & Quantum-resistant \\
Falcon-512       & 3   & No  & Quantum-resistant \\
X25519-ML-KEM-768 & 5  & No  & Hybrid (transitional) \\
\bottomrule
\end{tabular}
\end{table}

\paragraph{Time Horizon Risk ($\thr$)}
This factor operationalizes Mosca's inequality (Eq.~\ref{eq:mosca}). Let
$x$ denote the data retention period, $y$ the estimated migration time
based on the current migration phase, and $z = t_\text{CRQC} - t_\text{now}$
the time until a CRQC is available. The scoring function is:

\begin{equation}\label{eq:thr}
  \thr(a) =
  \begin{cases}
    100 & \text{if } x + y - z > 5 \\
    60 + \frac{(x + y - z)}{5} \cdot 40 & \text{if } 0 < x + y - z \leq 5 \\
    \max(0,\; 60 - |x + y - z| \cdot 10) & \text{if } x + y - z \leq 0
  \end{cases}
\end{equation}

\noindent Post-quantum algorithms receive a fixed $\thr = 5$ regardless
of the timeline, reflecting their inherent resistance.

The estimated migration time $y$ is determined by the asset's current
phase in a seven-stage migration continuum:

\begin{table}[t]
\centering
\caption{Migration phase model.}
\label{tab:migration-phases}
\footnotesize
\begin{tabular}{@{}lcc@{}}
\toprule
\textbf{Phase} & \textbf{Est. Years ($y$)} & \textbf{MGS} \\
\midrule
\textsc{Not Started}        & 5   & 100 \\
\textsc{Planning}           & 4   & 80  \\
\textsc{Testing}            & 3   & 60  \\
\textsc{Hybrid Deployment}  & 1   & 30  \\
\textsc{PQC Primary}        & 0.5 & 10  \\
\textsc{PQC Only}           & 0   & 2   \\
\textsc{CNSA 2.0 Compliant} & 0   & 0   \\
\bottomrule
\end{tabular}
\end{table}

\paragraph{Data Sensitivity Score ($\dss$)}
Assets are classified on a five-level scale: Public ($\dss{=}10$),
Internal ($30$), Confidential ($60$), Secret ($85$), and Top Secret
($100$). This weighting ensures that highly classified data requiring
long-term confidentiality receives higher migration priority.

\paragraph{Migration Gap Score ($\mgs$)}
The MGS quantifies how far an asset is from completing its PQC
migration, as shown in Table~\ref{tab:migration-phases}. An asset that
has not started migration scores $\mgs = 100$, while a fully
CNSA~2.0-compliant asset scores $\mgs = 0$.

\paragraph{Exposure Surface ($\es$)}
The exposure surface factor accounts for the number of systems or
endpoints using a given cryptographic asset:

\begin{equation}\label{eq:es}
  \es(a) =
  \begin{cases}
    10  & \text{if } n \leq 1 \\
    30  & \text{if } 1 < n \leq 10 \\
    60  & \text{if } 10 < n \leq 100 \\
    80  & \text{if } 100 < n \leq 1000 \\
    100 & \text{if } n > 1000
  \end{cases}
\end{equation}

\noindent where $n$ is the system count. A widely deployed algorithm
represents a larger attack surface and thus higher risk.

\subsection{Risk Classification}

The composite QRS maps to four risk levels:

\begin{equation}\label{eq:risk-levels}
  \text{Risk}(a) =
  \begin{cases}
    \textsc{Critical} & \text{if } \qrs(a) \in [76, 100] \\
    \textsc{High}     & \text{if } \qrs(a) \in [51, 75] \\
    \textsc{Moderate} & \text{if } \qrs(a) \in [26, 50] \\
    \textsc{Low}      & \text{if } \qrs(a) \in [0, 25]
  \end{cases}
\end{equation}

Each level triggers a prescribed action: \textsc{Critical} mandates
emergency migration within 90~days, \textsc{High} initiates a 6-month
migration plan, \textsc{Moderate} schedules the asset for the next
18-month cycle, and \textsc{Low} requires annual review only.

\subsection{Harvest-Now-Decrypt-Later Detection}

An asset $a$ is flagged for HNDL risk when its algorithm is vulnerable
to Shor's algorithm \emph{and} its data retention period exceeds the
time until a CRQC:

\begin{equation}\label{eq:hndl}
  \text{HNDL}(a) = \text{Shor-vulnerable}(a) \wedge (x > z)
\end{equation}

\noindent Assets flagged for HNDL risk receive immediate attention
regardless of their QRS, as the data they protect is already at risk
of future decryption.

\subsection{Portfolio Assessment}

The scoring engine aggregates individual assessments into a portfolio
view, computing: the mean QRS across all assets, the distribution of
risk levels, the percentage of assets at PQC or hybrid migration phases
(migration coverage), and the percentage meeting CNSA~2.0 requirements.
This enables executive-level reporting and resource allocation.

% =================================================================
%  V. HYBRID KEY EXCHANGE
% =================================================================
\section{Hybrid Key Exchange}\label{sec:hybrid}

During the transition period, hybrid key exchange constructions provide
defense-in-depth: if either the classical or the post-quantum component
is broken, the overall key exchange remains secure~\cite{giacon2018,
dowling2020}. We specify three variants targeting different deployment
contexts, following the IETF hybrid design~\cite{draft-hybrid-design}.

\subsection{Construction}

Let $\text{ECDH} = (\text{KeyGen}_C, \text{DH})$ be a classical
Diffie--Hellman key agreement and $\text{ML-KEM} = (\text{KeyGen}_P,
\text{Encaps}, \text{Decaps})$ be the post-quantum KEM. The hybrid
scheme $\text{HybridKEM}$ proceeds as follows:

\vspace{0.5em}
\noindent\textbf{Key Generation:}
\begin{align}
  (pk_C, sk_C) &\leftarrow \text{KeyGen}_C() \\
  (pk_P, sk_P) &\leftarrow \text{KeyGen}_P() \\
  pk &= pk_C \| pk_P,\quad sk = sk_C \| sk_P
\end{align}

\noindent\textbf{Encapsulation} (client):
\begin{align}
  (ct_C, K_C) &\leftarrow \text{ECDH}(pk_C) \\
  (ct_P, K_P) &\leftarrow \text{Encaps}(pk_P) \\
  ct &= ct_C \| ct_P \\
  K &= \text{HKDF}(K_C \| K_P,\; \text{salt},\; \text{info})
  \label{eq:combine}
\end{align}

\noindent\textbf{Decapsulation} (server):
\begin{align}
  K_C &\leftarrow \text{DH}(ct_C, sk_C) \\
  K_P &\leftarrow \text{Decaps}(ct_P, sk_P) \\
  K &= \text{HKDF}(K_C \| K_P,\; \text{salt},\; \text{info})
\end{align}

The shared-secret combination (Eq.~\ref{eq:combine}) uses
HKDF-SHA256~\cite{rfc5869} with:
\begin{itemize}[leftmargin=*]
\item $\text{salt} = \texttt{"qbitel-hybrid-kex-v1"}$ (fixed, public);
\item $\text{info} = \text{variant name}$ (e.g.,
  \texttt{"x25519-mlkem-768"}) for domain separation;
\item output length = 32 bytes (256 bits).
\end{itemize}

The concatenation of $K_C \| K_P$ as input key material (IKM) to HKDF
ensures that the derived key depends on both components. The domain
separation via the \texttt{info} parameter prevents cross-variant key
confusion.

\subsection{Variants}

Table~\ref{tab:hybrid-variants} summarizes the three variants.

\begin{table}[t]
\centering
\caption{Hybrid KEM variant specifications.}
\label{tab:hybrid-variants}
\footnotesize
\begin{tabular}{@{}lcccc@{}}
\toprule
\textbf{Variant} & \textbf{Classical} & \textbf{PQC} & \textbf{PK (B)} & \textbf{CT (B)} \\
\midrule
X25519-ML-KEM-768  & X25519 & ML-KEM-768  & 1\,216 & 1\,120 \\
P384-ML-KEM-1024   & P-384  & ML-KEM-1024 & 1\,665 & 1\,665 \\
X25519-ML-KEM-512  & X25519 & ML-KEM-512  & 832    & 800    \\
\bottomrule
\end{tabular}
\end{table}

\begin{itemize}[leftmargin=*]
\item \textbf{X25519-ML-KEM-768} is the recommended default for TLS~1.3,
  matching the deployment choice of Chrome and Cloudflare~\cite{cloudflare-pq,
  chrome-pq}. It provides NIST Level~3 security with moderate overhead
  (1\,216-byte public key, 1\,120-byte ciphertext).

\item \textbf{P384-ML-KEM-1024} targets enterprise and government
  environments requiring NIST Level~5 security. The use of P-384
  provides 192-bit classical security and compliance with CNSA~2.0
  requirements for the classical component.

\item \textbf{X25519-ML-KEM-512} is optimized for constrained
  environments (IoT, embedded systems) where bandwidth is limited. It
  provides NIST Level~1 security with the smallest footprint (832-byte
  public key, 800-byte ciphertext).
\end{itemize}

\subsection{Security Proof Sketch}

\begin{theorem}[Hybrid IND-CCA2 Security]
If either the classical ECDH component satisfies the decisional
Diffie--Hellman (DDH) assumption or the ML-KEM component is IND-CCA2
secure under Module-LWE, then the hybrid construction is IND-CCA2
secure, provided HKDF is modeled as a random oracle.
\end{theorem}

\begin{proof}[Proof sketch]
We proceed via a game-hopping argument following Giacon et al.~\cite{giacon2018}.

\emph{Game~0:} The real IND-CCA2 game. The challenger generates a hybrid
keypair and the adversary receives either the real or a random shared
secret.

\emph{Game~1:} Replace $K_C$ (the classical shared secret) with a
uniformly random value. The transition is justified by the DDH
assumption in the classical group.

\emph{Game~2:} Replace $K_P$ (the ML-KEM shared secret) with a uniformly
random value. The transition is justified by IND-CCA2 security of ML-KEM
under Module-LWE~\cite{regev2005, lyubashevsky2010}.

\emph{Game~3:} In Game~2, the IKM to HKDF is two independent random
values $K_C \| K_P$. Model HKDF as a random oracle; then the output $K$
is indistinguishable from random.

The adversary's advantage is bounded by:
\begin{equation}
  \text{Adv}^{\text{IND-CCA2}}_{\text{Hybrid}} \leq \text{Adv}^{\text{DDH}}_C + \text{Adv}^{\text{IND-CCA2}}_{\text{ML-KEM}} + \text{Adv}^{\text{PRF}}_{\text{HKDF}}
\end{equation}

Crucially, if either $\text{Adv}^{\text{DDH}}_C$ or
$\text{Adv}^{\text{IND-CCA2}}_{\text{ML-KEM}}$ is negligible, the
overall advantage remains negligible. This is the \emph{hybrid
defense-in-depth} property.
\end{proof}

\subsection{TLS 1.3 Integration}

The hybrid construction integrates into TLS~1.3~\cite{rfc8446} by
encoding the hybrid public key and ciphertext as opaque extensions in the
\texttt{key\_share} extension of ClientHello and ServerHello messages,
following the approach of draft-ietf-tls-hybrid-design~\cite{draft-hybrid-design}.
The hybrid shared secret replaces the standard ECDHE shared secret in the
TLS~1.3 key schedule; all subsequent key derivation (handshake keys,
application keys) proceeds identically.

\subsection{Memory Safety}

Private key material and shared secrets are zeroized upon deallocation.
The \texttt{HybridPrivateKey} and \texttt{HybridSharedSecret} objects
overwrite their byte arrays in the destructor, providing defense against
cold-boot and memory-scraping attacks. While Python's garbage collector
does not guarantee immediate collection, the explicit zeroization in
\texttt{\_\_del\_\_} reduces the exposure window.

% =================================================================
%  VI. IMPLEMENTATION AND EVALUATION
% =================================================================
\section{Implementation and Evaluation}\label{sec:eval}

We implement \system{} in Python using the \texttt{cryptography} library
for classical primitives and a custom ML-KEM implementation following
FIPS~203. The implementation comprises approximately 2\,400 lines of
code across three modules: \texttt{agility.py} (negotiation),
\texttt{quantum\_threat\_scoring.py} (scoring), and \texttt{hybrid.py}
(key exchange). All modules include Prometheus metrics instrumentation
for production monitoring.

\subsection{Experimental Setup}

Benchmarks were conducted on an Intel Xeon E5-2680 v4 (2.4\,GHz,
14 cores) with 64\,GB RAM running Ubuntu 22.04 and Python~3.11.
All timing results are averaged over 10\,000 iterations with 1\,000
warm-up iterations excluded. Cryptographic operations use
\texttt{cryptography}~42.0 with OpenSSL~3.2.

\subsection{Negotiation Performance}

Table~\ref{tab:negotiation-perf} reports negotiation latency for
different policy configurations.

\begin{table}[t]
\centering
\caption{Negotiation latency by policy (microseconds).}
\label{tab:negotiation-perf}
\footnotesize
\begin{tabular}{@{}lrrr@{}}
\toprule
\textbf{Policy} & \textbf{Mean ($\mu$s)} & \textbf{P99 ($\mu$s)} & \textbf{Algorithms} \\
\midrule
Legacy Compatible  & 145 & 312  & 28 \\
Transitional       & 187 & 398  & 22 \\
Post-Quantum       & 162 & 341  & 14 \\
CNSA 2.0           & 128 & 267  & 8  \\
\midrule
Renegotiation      & 203 & 445  & -- \\
\bottomrule
\end{tabular}
\end{table}

Negotiation completes in under 200\,$\mu$s on average across all
policies. The CNSA~2.0 policy achieves the lowest latency because its
strict constraints reduce the candidate set to 8 algorithms, minimizing
the intersection computation. Renegotiation (banning an algorithm and
recomputing) adds approximately 15\% overhead.

\subsection{Quantum Threat Scoring Throughput}

The scoring engine was evaluated on synthetic portfolios of varying size.
Table~\ref{tab:scoring-throughput} reports throughput in assets per
second.

\begin{table}[t]
\centering
\caption{Scoring engine throughput.}
\label{tab:scoring-throughput}
\footnotesize
\begin{tabular}{@{}rrrr@{}}
\toprule
\textbf{Portfolio Size} & \textbf{Throughput (assets/s)} & \textbf{Mean QRS} & \textbf{Critical (\%)} \\
\midrule
100       & 58\,200  & 52.3 & 28 \\
1\,000    & 55\,800  & 51.8 & 27 \\
10\,000   & 53\,400  & 52.1 & 29 \\
100\,000  & 51\,200  & 52.0 & 28 \\
\bottomrule
\end{tabular}
\end{table}

Throughput remains above 50\,000 assets/second even at 100\,000 assets,
with near-linear scaling. The computation is CPU-bound and dominated by
the Mosca inequality evaluation in the THR factor. The scoring function
requires no cryptographic operations, enabling batch processing of
large portfolios.

\subsection{Hybrid KEM Performance}

Table~\ref{tab:hybrid-perf} compares the three hybrid variants against
standalone ML-KEM and classical ECDH.

\begin{table}[t]
\centering
\caption{Hybrid KEM performance benchmarks.}
\label{tab:hybrid-perf}
\footnotesize
\begin{tabular}{@{}lrrrrr@{}}
\toprule
\textbf{Scheme} & \textbf{KeyGen} & \textbf{Encaps} & \textbf{Decaps} & \textbf{PK} & \textbf{CT} \\
 & \textbf{(ms)} & \textbf{(ms)} & \textbf{(ms)} & \textbf{(B)} & \textbf{(B)} \\
\midrule
X25519 (baseline)      & 0.03 & 0.04 & 0.04 & 32    & 32    \\
ML-KEM-768             & 0.12 & 0.15 & 0.14 & 1\,184 & 1\,088 \\
ML-KEM-1024            & 0.18 & 0.22 & 0.20 & 1\,568 & 1\,568 \\
ML-KEM-512             & 0.08 & 0.10 & 0.09 & 800   & 768   \\
\midrule
X25519-ML-KEM-768      & 0.16 & 0.21 & 0.19 & 1\,216 & 1\,120 \\
P384-ML-KEM-1024       & 0.26 & 0.33 & 0.30 & 1\,665 & 1\,665 \\
X25519-ML-KEM-512      & 0.12 & 0.15 & 0.14 & 832   & 800   \\
\bottomrule
\end{tabular}
\end{table}

The hybrid overhead relative to standalone ML-KEM is 10--15\% for
X25519-based variants and 30--45\% for P384-based variants, primarily due
to the additional elliptic-curve operation. The HKDF combination adds
negligible overhead ($<0.01$\,ms). All hybrid variants complete key
exchange in under 1\,ms, well within the latency budget for TLS
handshakes~\cite{paquin2020}.

\subsection{Comparison with Existing Frameworks}

Table~\ref{tab:comparison} positions \system{} against existing PQC
migration tools.

\begin{table}[t]
\centering
\caption{Comparison with existing frameworks.}
\label{tab:comparison}
\footnotesize
\begin{tabular}{@{}p{2.2cm}ccccc@{}}
\toprule
\textbf{Feature} & \textbf{\system{}} & \textbf{OQS} & \textbf{CBOM} & \textbf{IBM} \\
\midrule
Agility negotiation   & \checkmark & --         & --         & Partial \\
Threat scoring        & \checkmark & --         & \checkmark & \checkmark \\
Hybrid KEM            & \checkmark & \checkmark & --         & \checkmark \\
Policy enforcement    & \checkmark & --         & --         & \checkmark \\
Zero-downtime rotate  & \checkmark & --         & --         & -- \\
Mosca integration     & \checkmark & --         & \checkmark & -- \\
CNSA 2.0 compliance   & \checkmark & Partial    & --         & \checkmark \\
TLS 1.3 integration   & \checkmark & \checkmark & --         & Partial \\
\bottomrule
\end{tabular}
\end{table}

Open Quantum Safe (OQS)~\cite{paquin2020} provides library-level
PQC support including hybrid TLS but lacks higher-level agility
and risk assessment. Cryptographic Bill of Materials (CBOM)
approaches~\cite{dhs-pqc-roadmap} focus on inventory and
risk scoring but do not provide runtime negotiation or hybrid
constructions. IBM's quantum-safe toolkit offers threat assessment and
hybrid support but lacks the protocol-level agility and zero-downtime
rotation capabilities of \system{}.

% =================================================================
%  VII. SECURITY ANALYSIS
% =================================================================
\section{Security Analysis}\label{sec:security}

\subsection{Negotiation Security}

\begin{property}[Downgrade Resistance]
The negotiation protocol resists downgrade attacks: a man-in-the-middle
cannot force peers to negotiate a weaker algorithm than the strongest
mutually supported option, provided both peers enforce their policies.
\end{property}

This follows from the initiator-preference ordering: the protocol
selects the first algorithm in the initiator's preference list that
appears in the responder's capabilities. An attacker cannot remove
algorithms from the capability advertisement without detection, as the
advertisement is integrity-protected within the enclosing protocol
(e.g., TLS handshake).

\begin{property}[Policy Completeness]
For any security era $e$ and minimum security level $\ell$, the policy
engine correctly partitions the algorithm registry into allowed and
disallowed sets.
\end{property}

The filtering function applies four conjunctive checks: (i)~algorithm
not banned, (ii)~security bits $\geq \ell$, (iii)~algorithm status
$\neq \textsc{Banned}$, and (iv)~era compatibility. An algorithm is
admitted only if all four checks pass.

\subsection{Scoring Model Validity}

The QRS model is calibrated against two external references:

\begin{enumerate}[leftmargin=*]
\item \textbf{NIST SP 800-227 Timeline:} The default CRQC timeline
  ($t_\text{CRQC} = 2035$) aligns with NIST's planning horizon.
\item \textbf{CNSA 2.0 Deadlines:} The $\mgs$ factor correctly
  assigns $\mgs = 0$ to CNSA~2.0-compliant assets and $\mgs = 100$
  to assets that have not begun migration.
\end{enumerate}

The weighting vector $\mathbf{w} = (0.30, 0.25, 0.20, 0.15, 0.10)$
reflects the consensus view that algorithm vulnerability is the primary
driver of quantum risk, followed by the time horizon. We acknowledge
that the optimal weighting may vary by sector and regulatory context;
the framework supports custom weight configuration.

\subsection{Hybrid KEM Security}

The security of the hybrid construction rests on the KEM combiner
framework of Giacon et al.~\cite{giacon2018}: the concatenation-then-KDF
combiner preserves IND-CCA2 security if at least one component is
IND-CCA2 secure and the KDF is modeled as a random oracle. Our use of
HKDF-SHA256 with domain separation via the \texttt{info} parameter
satisfies this requirement under the PRF-ODH
assumption~\cite{brendel2022}.

\subsection{Limitations}

We identify the following limitations:

\begin{enumerate}[leftmargin=*]
\item \textbf{Python performance ceiling:} While sufficient for the
  control plane (negotiation, scoring), a production data-plane
  implementation would benefit from native code for the hybrid KEM.

\item \textbf{Side-channel resistance:} The current implementation does
  not provide constant-time guarantees for all operations. The ML-KEM
  encapsulation/decapsulation operations rely on the underlying library's
  constant-time properties.

\item \textbf{Memory zeroization:} Python's garbage collector does not
  guarantee immediate deallocation; the \texttt{\_\_del\_\_} zeroization
  is best-effort.

\item \textbf{Scoring model subjectivity:} The QRS weights and
  thresholds reflect expert judgment and may require sector-specific
  calibration.
\end{enumerate}

% =================================================================
%  VIII. RELATED WORK
% =================================================================
\section{Related Work}\label{sec:related}

\paragraph{NIST PQC Migration Guidance}
NIST SP~800-227~\cite{nist-sp800-227} provides a strategic roadmap but
does not specify runtime negotiation mechanisms or quantitative risk
scoring. Our agility protocol and QRS model operationalize the NIST
recommendations into deployable components.

\paragraph{Crypto Agility}
Ott and Peikert~\cite{ott2019} identified crypto agility as a key
research challenge. Housley~\cite{housley2024} explored algorithm agility
for digital certificates. DHS CISA~\cite{dhs-pqc-roadmap} published a
migration roadmap emphasizing inventory and planning phases. Our work
extends these efforts with a formal negotiation protocol and
policy-driven automation.

\paragraph{Quantum Risk Assessment}
Joseph et al.~\cite{joseph2022} provided a framework for organizational
PQC transition planning. Mosca~\cite{mosca2018} formalized the migration
urgency criterion. Our QRS model synthesizes these ideas into a
computable, multi-factor risk score suitable for portfolio-wide
assessment.

\paragraph{Hybrid Key Exchange}
The IETF hybrid design draft~\cite{draft-hybrid-design} specifies
hybrid key exchange for TLS. Stebila et al.~\cite{bos2018} and
Paquin et al.~\cite{paquin2020} benchmarked post-quantum TLS.
Giacon et al.~\cite{giacon2018} formalized KEM combiners, and
Dowling et al.~\cite{dowling2020} analyzed hybrid authenticated key
exchange. Schwabe et al.~\cite{schwabe2024} explored KEMTLS with
pre-distributed keys. Cloudflare~\cite{cloudflare-pq} and
Chrome~\cite{chrome-pq} deployed X25519-ML-KEM-768 in production.
Our hybrid constructions follow the concatenation-then-KDF approach
with explicit domain separation and extend it with enterprise
(P384-ML-KEM-1024) and constrained (X25519-ML-KEM-512) variants.

\paragraph{Open Quantum Safe}
The OQS project~\cite{paquin2020} provides a C library and OpenSSL
integration for PQC algorithms. While OQS excels at library-level
support, it does not address organizational policy enforcement,
risk-based migration prioritization, or runtime algorithm negotiation.
\system{} is complementary to OQS and can use it as a backend for
cryptographic operations.

% =================================================================
%  IX. CONCLUSION
% =================================================================
\section{Conclusion}\label{sec:conclusion}

We have presented \system{}, an integrated framework for managing the
post-quantum cryptographic transition. The Crypto Agility Negotiation
Protocol enables runtime algorithm selection with policy enforcement and
zero-downtime rotation. The Quantum Threat Scoring Engine provides
quantitative risk prioritization via a five-factor composite model
calibrated against NIST and CNSA~2.0 timelines. Three hybrid key exchange
constructions offer defense-in-depth for TLS~1.3 with practical
performance characteristics.

Our evaluation demonstrates that the framework is suitable for
large-scale deployment: negotiation latency remains under 200\,$\mu$s,
the scoring engine processes over 50\,000 assets per second, and hybrid
KEM overhead is within 15\% of standalone ML-KEM. Future work includes
native-code acceleration for the hybrid KEM data plane, formal
verification of the negotiation protocol, and sector-specific QRS
weight calibration for financial services, healthcare, and critical
infrastructure.

% =================================================================
%  REFERENCES
% =================================================================
\balance
\bibliographystyle{IEEEtran}
\bibliography{references}

\end{document}
